\documentclass[12pt,a4paper]{article}

\usepackage[margin=2.5cm]{geometry}
\usepackage{booktabs}
\usepackage{array}
\usepackage{xcolor}
\usepackage{hyperref}
\usepackage{enumitem}
\usepackage{titlesec}
\usepackage{parskip}
\usepackage{fancyhdr}
\usepackage{lmodern}
\usepackage[T1]{fontenc}

\definecolor{springerblue}{RGB}{26,58,107}

\hypersetup{
  colorlinks=true,
  linkcolor=springerblue,
  urlcolor=springerblue
}

\titleformat{\section}
  {\large\bfseries\color{springerblue}}
  {\thesection.}{0.6em}{}[\vspace{-4pt}\rule{\linewidth}{0.4pt}]

\titleformat{\subsection}
  {\normalsize\bfseries}
  {}{0em}{}

\pagestyle{fancy}
\fancyhf{}
\fancyhead[L]{\small\textcolor{springerblue}{Contribution Information Sheet}}
\fancyhead[R]{\small\textcolor{springerblue}{Discover Artificial Intelligence}}
\fancyfoot[C]{\thepage}
\renewcommand{\headrulewidth}{0.4pt}

% ---------------------------------------------------------------
\begin{document}

% Title block
\begin{center}
  {\LARGE\bfseries\color{springerblue} Contribution Information Sheet}\\[6pt]
  {\large Discover Artificial Intelligence --- Springer Nature}\\[18pt]
  \rule{\linewidth}{1pt}\\[10pt]
  {\large\bfseries Optimised Support Vector Regression for California Housing\\
  Price Prediction: The Critical Role of Feature Engineering\\
  and Hyperparameter Tuning}\\[10pt]
  \rule{\linewidth}{0.4pt}
\end{center}

\vspace{6pt}

\begin{tabular}{@{}ll@{}}
\textbf{Corresponding Author:}  & Emmanuel Adutwum \\
\textbf{Affiliation:}           & Soka University of America, Aliso Viejo, CA, USA \\
\textbf{Email:}                 & \href{mailto:eadutwum@soka.edu}{eadutwum@soka.edu} \\
\textbf{Transfer Ref.\ (SNAPP):}& 80b45e30-56da-4033-b5e6-07f0e94c6143 \\
\end{tabular}

\vspace{14pt}

% ---------------------------------------------------------------
\section{Scientific Question Addressed}

A recent comparative study (Preethi et al., \textit{SN Computer Science}, 2025) reported
that Support Vector Regression (SVR) ranks \emph{last} among five regression algorithms on
the California Housing benchmark, achieving R\textsuperscript{2} = 0.60, and implicitly
concluded that SVR is a weak choice for this class of problem.

This paper asks a more precise question: \textbf{Does that poor performance reflect an
inherent limitation of SVR as an algorithm, or does it reflect the experimental
configuration under which SVR was evaluated?}

This is a question about \textbf{evaluation methodology and reproducibility in applied AI}
--- specifically, how much predictive performance is left unrealised when standard
preprocessing steps are omitted from kernel-based models, and whether algorithm-ranking
conclusions drawn from under-configured experiments are valid.

% ---------------------------------------------------------------
\section{Methodology}

A four-stage ablation workflow is applied entirely within a leakage-safe
\texttt{scikit-learn} Pipeline:

\begin{enumerate}[leftmargin=1.8em]
  \item \textbf{Baseline:} SVR-RBF on raw, unscaled features --- replicating Preethi
        et al.\ (2025).
  \item \textbf{Scaling:} \texttt{StandardScaler} applied inside the Pipeline so scaling
        statistics are recomputed on each cross-validation fold (no data leakage).
  \item \textbf{Feature Engineering:} Ten domain-motivated features derived from the
        eight raw inputs (rooms-per-household, bedrooms ratio, population density,
        log-median income, geographic interaction terms). An ensemble importance score
        combining Mutual Information (40\%), Pearson correlation (30\%), and Random
        Forest importance (30\%) selects the top 12 features.
  \item \textbf{Hyperparameter Tuning:} \texttt{RandomizedSearchCV} with 20 iterations
        and 3-fold inner cross-validation identifies the optimal kernel, $C$,
        $\varepsilon$, and $\gamma$.
\end{enumerate}

A formal ablation study isolates the R\textsuperscript{2} contribution of each stage.
Ten-fold outer cross-validation provides confidence intervals on the final generalisation
estimate.

% ---------------------------------------------------------------
\section{Key Findings}

\vspace{4pt}
\begin{center}
\begin{tabular}{lcc}
\toprule
\textbf{Pipeline Stage} & \textbf{Test R\textsuperscript{2}} & \textbf{Gain} \\
\midrule
Baseline --- no scaling (replication of prior work) & $-$0.054 & --- \\
\quad + Feature scaling only                        & 0.690    & \textbf{+0.744} \\
\quad + Domain-motivated feature engineering        & 0.716    & +0.026 \\
\quad + Hyperparameter tuning (RandomizedSearchCV)  & \textbf{0.723} & +0.008 \\
\bottomrule
\end{tabular}
\end{center}
\vspace{6pt}

\textbf{Primary finding:} Feature scaling alone accounts for 95.7\% of the total
R\textsuperscript{2} improvement (+0.744), reversing SVR from negative predictive
power to competitive performance. This is not a minor artefact.

\textbf{Secondary finding:} The properly configured SVR-RBF ranks 4th in a 10-model
comparison --- below XGBoost (0.832), Random Forest (0.814), and Gradient Boosting
(0.783), but substantially above Decision Trees, Linear Regression, Ridge, Lasso,
Elastic Net, and KNN.

\textbf{Generalisation:} Ten-fold cross-validation yields mean R\textsuperscript{2} =
0.703 (95\% CI: [0.630, 0.775]). The confidence interval does not overlap with the
reported baseline of 0.60, confirming statistical robustness across splits.

% ---------------------------------------------------------------
\section{Specific Contributions to the AI Field}

\begin{enumerate}[leftmargin=1.8em]
  \item \textbf{Ablation framework for pipeline evaluation.} A replicable four-stage
        ablation design for isolating the contribution of scaling, feature engineering,
        and hyperparameter optimisation. Applicable to any kernel-based or
        distance-sensitive algorithm.

  \item \textbf{Quantification of preprocessing sensitivity.} The +0.744 R\textsuperscript{2}
        gain from scaling alone on a widely-used benchmark provides a concrete,
        citable data point corroborating Fernandez-Delgado et al.\ (2014) and
        Cao \& Tay (2001) on a modern standardised benchmark.

  \item \textbf{Leakage-safe reproducible implementation.} The \texttt{scikit-learn}
        Pipeline with \texttt{ColumnTransformer} ensures all preprocessing statistics
        are computed exclusively on training data within each fold. Full code is
        released as open source.

  \item \textbf{Evaluation practice.} The paper demonstrates that algorithm rankings
        derived from under-configured experiments can be misleading --- with direct
        implications for comparative benchmarking standards in the AI/ML community.
\end{enumerate}

% ---------------------------------------------------------------
\section{Significance and Fit for Discover Artificial Intelligence}

The scope of \textit{Discover Artificial Intelligence} explicitly covers
\emph{``machine and deep learning,''} \emph{``data analytics,''} and
\emph{``AI enabled data-driven techniques''} --- all directly addressed here.

Beyond the housing application, the paper raises a question fundamental to applied AI:
\textbf{how much of an algorithm's perceived performance is a property of the algorithm
itself, and how much is a property of the experimental pipeline around it?} The answer
--- that 95.7\% of SVR's improvement comes from a single preprocessing step --- is
both striking and broadly generalisable to any practitioner evaluating kernel methods
against ensemble alternatives.

The California Housing dataset (20,640 instances, part of the \texttt{scikit-learn}
standard library) is one of the most widely cited supervised learning benchmarks,
making results immediately reproducible by any reader with a Python environment.

% ---------------------------------------------------------------
\section{Intended Audience}

\begin{itemize}[leftmargin=1.8em]
  \item ML practitioners evaluating SVR and kernel methods for regression tasks
  \item Researchers designing comparative benchmarking studies
  \item AutoML and pipeline design researchers (preprocessing sensitivity is a
        core concern)
  \item Data science educators using the California Housing dataset in coursework
\end{itemize}

% ---------------------------------------------------------------
\section{Reproducibility}

Full implementation is publicly available at:\\[4pt]
\url{https://github.com/emmanueladutwum123/svr-california-housing}

\vspace{4pt}
\begin{tabular}{@{}ll@{}}
  \texttt{openwork.tex}          & LaTeX source \\
  \texttt{svr\_implementation.py}& Complete Python pipeline \\
  \texttt{svr\_notebook.ipynb}   & Jupyter notebook with all outputs \\
  \texttt{Picture1--13.png}      & All manuscript figures \\
\end{tabular}

\vspace{6pt}
Dataset: \texttt{sklearn.datasets.fetch\_california\_housing} (public, no special
access required). All results reproducible with scikit-learn $\geq$ 1.0.

% ---------------------------------------------------------------
\section{Keywords}

Support Vector Regression; Feature Engineering; Hyperparameter Optimisation;
Ablation Study; Preprocessing Sensitivity; California Housing Dataset;
Scikit-learn Pipeline; Kernel Methods; Reproducible Machine Learning;
Supervised Regression

% ---------------------------------------------------------------
\section{Conflict of Interest}

The author declares no conflicts of interest.

% ---------------------------------------------------------------
\section{Data Availability}

The California Housing dataset is publicly available via the scikit-learn Python
library (\texttt{sklearn.datasets.fetch\_california\_housing}), originally published
by Pace \& Barry (1997). No proprietary data are used. All code and experimental
outputs are available at
\url{https://github.com/emmanueladutwum123/svr-california-housing}.

\vspace{20pt}
\noindent\rule{\linewidth}{0.4pt}\\[4pt]
\small\textit{Prepared for submission to Discover Artificial Intelligence,
Springer Nature, 2026.}

\end{document}
